\section{August 24, 2026}

\subsection{Fields}

\begin{definition}[Binary operation]
  A \emph{binary operation} on a set $S$ is a function
  \[
    F\colon S\times S\longrightarrow S.
  \]
  Thus, for every $x,y\in S$, the value $F(x,y)$ is again an element of
  $S$. We usually write $x\cdot y$ or $x+y$ in place of $F(x,y)$.
\end{definition}

\begin{example}
  Addition on $\bb{Z}$ is a binary operation: the map
  \[
    \bb{Z}\times\bb{Z}\longrightarrow\bb{Z},
    \qquad (m,n)\longmapsto m+n
  \]
  takes every pair of integers to an integer.
\end{example}

\begin{definition}
  Let $F$ be a binary operation on $S$. A subset $T\subseteq S$ is
  \emph{closed under $F$} if
  \[
    F(x,y)\in T
  \]
  whenever $x,y\in T$.
\end{definition}

\begin{example}
  The positive integers $\bb{Z}_{>0}$ are closed under addition. The odd
  integers are not closed under addition, but they are closed under
  multiplication.
\end{example}

\begin{definition}[Field]
  Let $\bb{F}$ be a set equipped with two binary operations, called
  \emph{addition} and \emph{multiplication}. The sum and product of
  $a,b\in\bb{F}$ are denoted by $a+b$ and $ab$, respectively. Suppose that
  for every $a,b,c\in\bb{F}$,
  \begin{enumerate}[label=(\roman*)]
    \item $a+b=b+a$;
    \item $(a+b)+c=a+(b+c)$;
    \item $ab=ba$;
    \item $a(bc)=(ab)c$;
    \item $a(b+c)=ab+ac$.
  \end{enumerate}
  Suppose further that
  \begin{enumerate}[label=(\roman*),resume]
    \item $\bb{F}$ contains an additive identity $0$ and a multiplicative
      identity $1$, so that
      \[
        a+0=a \qquad\text{and}\qquad 1a=a
      \]
      for every $a\in\bb{F}$;
    \item $1\neq 0$;
    \item every $a\in\bb{F}$ has an additive inverse $-a\in\bb{F}$ such
      that $a+(-a)=0$;
    \item every nonzero $a\in\bb{F}$ has a multiplicative inverse
      $a^{-1}\in\bb{F}$ such that $aa^{-1}=1$.
  \end{enumerate}
  Then $\bb{F}$ is called a \emph{field}.
\end{definition}

\begin{example}
  The integers $\bb{Z}$ do not form a field: for example, $2$ has no
  multiplicative inverse in $\bb{Z}$. The sets $\bb{Q}$, $\bb{R}$, and
  $\bb{C}$, with their usual operations, are fields.
\end{example}

There are also finite fields. For example,
\[
  \bb{F}_2=\{0,1\}
  \qquad\text{and}\qquad
  \bb{F}_3=\{0,1,2\},
\]
where the elements are the possible remainders upon division by $2$ and
$3$, respectively. Addition and multiplication in $\bb{F}_3$ are given by
the following tables:
\[
\begin{array}{c|ccc}
  + & 0 & 1 & 2 \\
  \hline
  0 & 0 & 1 & 2 \\
  1 & 1 & 2 & 0 \\
  2 & 2 & 0 & 1
\end{array}
\qquad\qquad
\begin{array}{c|ccc}
  \cdot & 0 & 1 & 2 \\
  \hline
  0 & 0 & 0 & 0 \\
  1 & 0 & 1 & 2 \\
  2 & 0 & 2 & 1
\end{array}.
\]
More generally, if $p$ is prime, then
\[
  \bb{F}_p=\{0,1,\ldots,p-1\},
\]
with addition and multiplication performed modulo $p$, is a field.

\begin{lemma}
  If $a$ is an element of a field $\bb{F}$, then $a\cdot 0=0$.
\end{lemma}

\begin{proof}
  Since $0+0=0$, distributivity gives
  \[
    a\cdot 0=a(0+0)=a\cdot 0+a\cdot 0.
  \]
  Adding the additive inverse of $a\cdot 0$ to both sides yields
  $0=a\cdot 0$.
\end{proof}
